\documentclass[11pt]{article}

\usepackage[T1]{fontenc}
\usepackage[utf8]{inputenc}
\usepackage{mathpazo}
\usepackage[scaled=0.90]{helvet}
\usepackage{amsmath,amssymb}
\usepackage{microtype}
\usepackage[a4paper,top=2.2cm,bottom=2.2cm,left=2.4cm,right=2.4cm]{geometry}
\usepackage{parskip}
\usepackage{enumitem}
\setlist[itemize]{leftmargin=1.4em,itemsep=2pt,topsep=3pt}
\setlist[enumerate]{leftmargin=1.6em,itemsep=2pt,topsep=3pt}
\usepackage{xcolor}
\definecolor{accent}{RGB}{20,60,110}
\usepackage[hidelinks]{hyperref}
\usepackage{titlesec}
\titleformat{\section}{\normalfont\sffamily\bfseries\large\color{accent}}{\thesection.}{0.6em}{}
\titleformat{\subsection}{\normalfont\sffamily\bfseries\normalsize\color{accent}}{\thesubsection}{0.6em}{}
\titlespacing*{\section}{0pt}{12pt}{4pt}
\titlespacing*{\subsection}{0pt}{8pt}{2pt}
\usepackage{fancyhdr}
\pagestyle{fancy}
\fancyhf{}
\renewcommand{\headrulewidth}{0pt}
\fancyfoot[C]{\footnotesize\sffamily\color{gray}\thepage}

\begin{document}

\begin{center}
{\footnotesize\sffamily\color{gray}
 SAPIENZA UNIVERSITY OF ROME\\[1pt]
 Master's Degree in Computer Engineering \textbullet\
 Excellence Path (\textit{Percorso di Eccellenza})}\\[10pt]
{\LARGE\sffamily\bfseries\color{accent}
 Anonymous Credentials in Practice:}\\[3pt]
{\large\sffamily\bfseries\color{accent}
 A Study of Device Binding, Zero-Knowledge Proofs,\\
 and the Road to Standardisation}\\[10pt]
{\normalsize\sffamily
 \textbf{Author:} Federico Rufini \qquad
 \textbf{Supervisor:} Prof.\ Visconti}\\[2pt]
{\footnotesize\sffamily\color{gray} Academic Year 2025/2026}
\end{center}

\vspace{2pt}
{\color{accent}\hrule height 0.9pt}
\vspace{8pt}

\noindent\textit{\small
Anonymous credential systems promise to let users prove attributes about
themselves---age, nationality, professional qualifications---without
revealing their identity or allowing different presentations to be linked.
This report documents the work carried out during my Excellence Path, which
consisted of a focused study of the recent scientific literature on
privacy-preserving credentials, with particular attention to two influential
lines of research: anonymous credentials built on top of the ECDSA signature
scheme as found in deployed secure hardware (Frigo and Shelat, 2024), and
device-bound anonymous credentials designed to remain private even under
partial or total hardware compromise (Friedrichs et al., 2025). The study
is framed by the imminent rollout of the EU Digital Identity (EUDI) Wallet,
which turns these formerly academic questions into a matter of practical
urgency.}

\vspace{6pt}

\section{Motivation}

Anonymous credentials are one of the oldest studied objects in applied
cryptography, yet their deployment in real systems has remained elusive for
decades. The core tension is well known: to be trustworthy, a credential must
be tied to a device that cannot be cloned; but the same binding that prevents
cloning also enables tracking. Breaking this tension without modifying
existing hardware or relying on non-standard cryptographic assumptions is the
central open problem of the field.

Under the guidance of Prof.\ Visconti, I used my Excellence Path to study
this problem through its primary literature. The goal was not a broad survey
but a deep reading of a small number of foundational and recent papers,
sufficient to understand their technical contributions, to evaluate them
critically, and to locate the open problems they leave behind. The choice of
papers was guided by a combination of cryptographic depth and practical
relevance: both works studied are directly applicable to systems being
standardised and deployed today.

\section{Background: the hardware constraint}

An anonymous credential system allows an \emph{issuer} to sign a set of
attributes, and a \emph{holder} to later prove possession of a valid
signature---and statements about the signed attributes---to a
\emph{verifier}, without revealing the signature itself or allowing the
verifier to link two presentations. The classical construction by Camenisch
and Lysyanskaya achieves this through algebraic signatures over which
efficient zero-knowledge proofs of knowledge can be built. More recently,
\textbf{BBS signatures}---now standardised by the IETF---have emerged as
the preferred instantiation: they support multi-message signing, selective
disclosure, and efficient ZK proof protocols built on pairings.

Device binding adds a further requirement: the credential must be
cryptographically tied to a specific device so that it cannot be trivially
copied. The natural mechanism is to keep a secret key inside a
\emph{Secure Element} (SE), a tamper-resistant chip that performs
cryptographic operations without ever exporting the key. The problem is that
today's Secure Elements---those found in hundreds of millions of shipped
smartphones---support only \textbf{ECDSA over NIST P-256}. They cannot
natively compute BBS or BLS signatures, pairing operations, or
zero-knowledge proofs. Any scheme requiring the hardware to do more therefore
cannot be deployed without a new generation of devices. This constraint is
the common starting point of both papers studied, and the main reason the
problem remains open.

\section{Longfellow: anonymous credentials from ECDSA (Frigo and Shelat, 2024)}

\subsection{Contribution}

The paper asks a deceptively simple question: can one build a fully-featured
anonymous credential system using \emph{only} ECDSA, the primitive that
existing secure hardware already supports? The construction---called
\textbf{Longfellow}---answers yes, and is the most technically demanding
work studied during the path.

The key insight is that the Secure Element does not need to compute a
privacy-preserving signature. Instead, it signs a commitment using its
standard ECDSA key, and the holder produces a \emph{zero-knowledge proof}
establishing that (a)~there exists a valid ECDSA signature over the
commitment, (b)~the commitment opens to a credential issued by a trusted
issuer, and (c)~the credential contains the claimed attribute values. The
verifier learns nothing beyond the truth of these statements. Because nothing
about the ECDSA key or the credential is revealed, multiple presentations
are unlinkable.

\subsection{Proof system architecture}

The ZK proof system is a carefully engineered combination of components.
\textbf{Ligero} (MPC-in-the-head) serves as the core proof system: it
requires no trusted setup, relies only on collision-resistant hash functions
(making it post-quantum secure in isolation), and supports arbitrary
arithmetic circuits. To avoid the bottleneck of heavy number-theoretic
transforms, Ligero is integrated with a \textbf{sum-check} protocol, which
allows the verifier to check polynomial sums over boolean hypercubes without
NTT operations. Polynomial commitments are instantiated via
\textbf{Reed-Solomon codes}, using additive FFT over binary fields and
linear convolution over NIST P-256.

The most original technical contribution is the treatment of two
sub-computations that naturally live in different fields. ECDSA verification
works over a prime field; SHA-256 hashing works over a binary field
(GF$(2^{128})$). The paper introduces two separate optimised circuits, one
for each field, connected by a \textbf{MAC-based cross-field consistency
check}: the witness shared between the circuits is authenticated so that a
cheating prover cannot substitute inconsistent values across the field
boundary. The resulting proof is made non-interactive via the
\textbf{Fiat--Shamir} heuristic, producing a single message that any
verifier can check without further interaction---a hard requirement of
OpenID4VP and ISO~18013-5. The European fork \emph{Longfellow-zk} (Dyne.org)
subsequently added SIMD128 assembler optimisations targeting European
credential formats.

\subsection{Performance}

The reported figures are remarkable for a general ZK system: ECDSA proofs
in approximately 60\,ms, and a full ISO MDOC presentation proof in roughly
1.2\,seconds on mobile hardware depending on credential size. The work
demonstrates convincingly that ZK-based, hardware-compatible anonymous
credentials are no longer purely theoretical.

\subsection{Critical assessment}

Several limitations emerged from the study. First, the post-quantum picture
is mixed: Ligero's security rests only on hash functions and is post-quantum
safe, but the ECDSA signature from the Secure Element is broken by Shor's
algorithm. The scheme is therefore best understood as a transitional
solution, viable until Secure Elements migrate to post-quantum primitives.
Second, the SE remains a \emph{single point of failure}: a compromised SE key
collapses user privacy entirely, and the paper proposes no multi-party
technique to distribute this trust. Third, the codebase exceeds 20\,000
lines of code and the authors have already issued corrections for
implementation errors---a reminder that the gap between a sound proof system
and a secure implementation is not trivial to close.

\section{Device-Bound Anonymous Credentials With(out) Trusted Hardware\\
\small\normalfont Friedrichs, Lehmann, Sidorenko, and Zacharakis, 2025}

\subsection{Contribution}

This paper approaches device binding from a different direction. Rather than
optimising for compatibility with existing ECDSA hardware, it asks: what
privacy guarantees are achievable when the hardware itself cannot be fully
trusted? The motivation is grounded in two concrete scenarios already present
in the EUDI ecosystem:
\begin{itemize}
  \item Many smartphones lack a locally certified SE at the required Level
        of Assurance (LoA High), so the regulation foresees the use of
        \emph{Remote HSMs}---cloud-based security modules outside the user's
        physical control.
  \item Even local Secure Elements are ``black boxes'': subverted firmware
        could embed a subliminal channel that leaks the user's identity
        across presentations without any party being able to detect it.
\end{itemize}

The paper's first contribution is a formal hierarchy of four unlinkability
levels (Unlink-0 through Unlink-3) capturing progressively stronger
adversarial models, up to a SE that is fully under real-time adversary
control. This taxonomy is itself a contribution: it gives the field a
precise vocabulary for what ``device binding without trusted hardware''
actually means, something prior work had left informal.

\subsection{Three constructions}

All three proposed schemes use \textbf{BBS signatures} for the credential
itself and differ only in how the device contributes its binding signature:
\begin{enumerate}
  \item \textbf{BBS-Schnorr.} Device binding via a Schnorr signature.
        Efficient, but achieves only Unlink-1: a malicious SE can exploit
        the signature randomness to embed a subliminal channel.
  \item \textbf{BBS-BLS.} Replaces Schnorr with a deterministic BLS
        signature. Determinism eliminates subliminal channels (Unlink-2),
        but the SE still learns the message it signs, so presentation
        context leaks to the hardware.
  \item \textbf{BBS-BlindBLS.} The SE signs blindly, learning neither the
        message nor the presentation context. Achieves Unlink-3 and full
        \emph{SE-obliviousness}; the construction designed for Remote HSMs.
\end{enumerate}

\subsection{Critical assessment}

The work resolves two main weaknesses of Longfellow: it removes the SE
single point of failure by supporting remote hardware, and it replaces the
ECDSA circuit with BBS, enabling richer ZK proofs without hardware-specific
engineering. Limitations are of two kinds. Some are declared by the authors:
BBS-Schnorr is not subversion-resilient, and BBS-BLS is not SE-oblivious;
these are acknowledged trade-offs, not oversights. More fundamentally,
\emph{none of the three constructions is post-quantum secure}---BBS, BLS, and
Schnorr all rest on the hardness of discrete logarithm.

Two further limitations emerged from independent analysis. The Remote HSM
model introduces unavoidable \emph{metadata leakage through timing}: an
adversary observing network traffic can correlate HSM request timing with
user activity even without breaking any cryptographic primitive, a channel
the paper does not address. And permanent network connectivity is required,
making the scheme unsuitable for offline credential presentation---a real
requirement in practice, which the scheme can only partially mitigate by
pre-fetching non-revocation proofs when connectivity is available.

\section{Comparison and synthesis}

Reading the two papers together reveals a clear complementarity. Longfellow
optimises for compatibility with today's deployed hardware at the cost of a
complex, ECDSA-specific proof system. DBAC optimises for a richer privacy
model---including hardware subversion and Remote HSM scenarios---at the cost
of requiring devices that can compute BBS or BLS. Neither achieves
post-quantum security end-to-end, and neither fully resolves the tension
between hardware binding and user privacy under a fully malicious device.

A third dimension became visible during the study. The recent announcement
that \textbf{Android~17} will expose \textbf{ML-DSA} (the NIST lattice-based
signature standard) through the standard Android Keystore TEE interface opens,
for the first time, the prospect of a post-quantum primitive natively
available in mobile secure hardware. ZK proof systems for lattice-based
signatures are still substantially heavier than their elliptic-curve
counterparts, so the problem is not immediately solved---but the landscape is
shifting. The DBAC unlinkability hierarchy provides a ready-made formal
framework for evaluating new constructions that will emerge in this setting.

\section{Relevance to the EU Digital Identity Wallet}

The EU Digital Identity Wallet---currently in its pre-deployment phase,
governed by the Architecture and Reference Framework (ARF)---is the largest
digital identity infrastructure project ever attempted in Europe, and the
closest real-world instantiation of the problems studied in this literature.
Its selective-disclosure mechanism relies on salted-hash SD-JWT Verifiable
Credentials and ISO/IEC~18013-5 mDOC presentations: mechanisms that inherit
the granularity and linkability limitations the anonymous-credential
literature has documented for decades. Disclosing a full date-of-birth field
to prove ``I am over 18'', or allowing multiple presentations to be
correlated through device-bound attribute signatures, are not hypothetical
issues but design choices embedded in the current standard.

The ARF itself acknowledges that zero-knowledge proofs are the only viable
long-term mitigation, yet mandates none. Longfellow has attracted serious
attention as a near-term upgrade path precisely because it requires no
hardware changes; the Dyne.org fork targets European credential formats
specifically. DBAC directly addresses the regulation's provision for Remote
HSMs, providing a formal treatment of the privacy implications the regulatory
text leaves underspecified. The literature studied during this path is
therefore not only academically relevant---it is a close reading of the
cryptographic infrastructure that will shape European digital identity for
the foreseeable future.

\section{Outcome of the path}

The Excellence Path produced a detailed, critical understanding of two of
the most technically substantive recent contributions to device-bound
anonymous credentials. Beyond the content, the work developed two capacities
harder to acquire from coursework alone: reading a complex cryptographic
paper and identifying not only what it claims but what it implicitly assumes
or leaves open; and situating a research contribution within a moving
practical context---hardware announcements, standardisation timelines,
regulatory choices---that determines which open problems are actually worth
pursuing. I am grateful to Prof.\ Visconti for the direction and the
critical engagement that shaped this study.

\vspace{8pt}
{\color{accent}\hrule height 0.4pt}
\vspace{4pt}
{\footnotesize\noindent\textbf{\sffamily References}\\[3pt]
[1]~M.~Frigo, a.~shelat, \textit{Anonymous Credentials from ECDSA}
(Longfellow), Google, 2024.\\[1pt]
[2]~J.~Friedrichs, A.~Lehmann, A.~Sidorenko, A.~Zacharakis,
\textit{Device-Bound Anonymous Credentials With(out) Trusted Hardware}, 2025.\\[1pt]
[3]~D.~Boneh, X.~Boyen, H.~Shacham, \textit{Short Group Signatures} (BBS);
IETF draft-irtf-cfrg-bbs-signatures.\\[1pt]
[4]~S.~Ames, C.~Hazay, Y.~Ishai, M.~Venkitasubramaniam,
\textit{Ligero: Lightweight Sublinear Arguments Without a Trusted Setup}, CCS~2017.\\[1pt]
[5]~A.~Lehmann, A.~Sidorenko, A.~Zacharakis,
\textit{Vision: A Modular Framework for Anonymous Credential Systems}.\\[1pt]
[6]~European Commission, \textit{Architecture and Reference Framework (ARF)
for the EUDI Wallet}, eudi.dev.\\[1pt]
[7]~OpenID Foundation, \textit{OpenID for Verifiable Presentations (OpenID4VP)};
ISO/IEC~18013-5:2021, \textit{Mobile driving licence (mDL) application}.\\[1pt]
[8]~Google Security Blog, \textit{Post-Quantum Cryptography in Android}
(ML-DSA, Android~17), 2026.
}

\end{document}